\chapter{Developer Experience}

\section{The Glue Philosophy}

Every powerful platform risks becoming inaccessible through its own complexity.
Connector solves this with \textbf{Glue} — a unified developer surface that
exposes the full power of the kernel, execution engines, and protocol stack
through a simple, human-readable grammar.

Glue is not a wrapper or an SDK. It is a \emph{language}: a set of verbs,
nouns, selectors, and time flags that compose into commands with consistent,
parseable output contracts. The same grammar works in the CLI, the REST API,
and the Rust library — so operators, developers, and automated pipelines all
speak the same language.

\begin{conceptbox}{Three design principles of Glue}
\begin{enumerate}[noitemsep]
  \item \textbf{Verbs first} — Commands read like instructions: \texttt{start agent}, \texttt{list agents}, \texttt{prove agent} — not \texttt{agent-start} or \texttt{agent.start()}.
  \item \textbf{Hide complexity, not power} — One verb can orchestrate many kernel operations. The complexity is hidden, but not removed — it is accessible through flags.
  \item \textbf{Instant value} — Every command produces useful output immediately, even with no flags. No ceremony, no configuration required to get started.
\end{enumerate}
\end{conceptbox}

\section{The Glue Grammar}

Glue commands follow a consistent four-part structure:

\begin{figure}[H]
\centering
\begin{tikzpicture}[
  part/.style={rectangle, rounded corners=3pt, draw=#1!70, fill=#1!10,
               text width=2.6cm, align=center, minimum height=1.1cm,
               font=\small\bfseries, inner sep=5pt},
  opt/.style={rectangle, rounded corners=3pt, draw=#1!40, fill=#1!5,
              text width=2.6cm, align=center, minimum height=1.1cm,
              font=\small\bfseries, inner sep=5pt, dashed},
  lbl/.style={font=\footnotesize\itshape, text=gray, align=center},
  node distance=0.6cm
]
\node[part=CKernel] (v) {\texttt{<verb>}};
\node[part=CExec, right=of v]   (n) {\texttt{<noun>}};
\node[part=CProto, right=of n]  (s) {\texttt{[selector]}};
\node[opt=CDX,   right=of s]    (f) {\texttt{[flags]}};

\node[lbl, below=0.25cm of v] {action};
\node[lbl, below=0.25cm of n] {resource type};
\node[lbl, below=0.25cm of s] {identity / id};
\node[lbl, below=0.25cm of f] {time, output, mode};

\node[above=0.3cm of v, font=\footnotesize, text=MidnightBlue]
  {\texttt{connectorctl}};

\end{tikzpicture}
\caption{Glue command grammar structure}
\label{fig:glue-grammar}
\end{figure}

\subsection{Verbs}

Glue provides 12 primary verbs that cover the full agent lifecycle:

\begin{table}[H]
\centering
\renewcommand{\arraystretch}{1.4}
\rowcolors{2}{gray!5}{white}
\begin{tabular}{@{}llp{7cm}@{}}
\toprule
\textbf{Verb} & \textbf{Category} & \textbf{What it does} \\
\midrule
\texttt{start}   & Lifecycle   & Start an agent, service, or node \\
\texttt{stop}    & Lifecycle   & Stop an agent or service gracefully \\
\texttt{restart} & Lifecycle   & Stop then start, preserving checkpoints \\
\texttt{list}    & Query       & List agents, sessions, memory, contracts \\
\texttt{show}    & Query       & Show summary state for one resource \\
\texttt{inspect} & Query       & Deep inspection — forensic mode available \\
\texttt{trace}   & Diagnostics & Decision trace with time-window filter \\
\texttt{prove}   & Evidence    & Generate cryptographic proof for an agent \\
\texttt{explain} & Evidence    & Decision-first explanation (\textit{why} an agent acted) \\
\texttt{deploy}  & Operations  & Deploy an agent manifest or contract \\
\texttt{backup}  & Maintenance & Checkpoint memory to an archive \\
\texttt{restore} & Maintenance & Restore from a backup or checkpoint \\
\bottomrule
\end{tabular}
\caption{Glue's 12 primary verbs}
\end{table}

\subsection{Nouns}

Ten primary nouns represent the resources Glue operates on:
\texttt{agent}, \texttt{session}, \texttt{memory}, \texttt{contract},
\texttt{tool}, \texttt{port}, \texttt{node}, \texttt{cluster},
\texttt{key}, \texttt{user}.

\subsection{Time Flags}

Time flags enable point-in-time and range queries against the immutable
MemPacket history — a direct expression of the kernel's content-addressed store:

\begin{table}[H]
\centering
\renewcommand{\arraystretch}{1.35}
\begin{tabular}{@{}lll@{}}
\toprule
\textbf{Flag} & \textbf{Format} & \textbf{Example} \\
\midrule
\texttt{--at}      & ISO 8601 or Unix timestamp & \texttt{--at 2026-03-22T14:00:00Z} \\
\texttt{--since}   & ISO 8601 or relative       & \texttt{--since 1h} \\
\texttt{--last}    & Duration shorthand         & \texttt{--last 15m} \\
\texttt{--before}  & Timestamp or event name    & \texttt{--before incident-001} \\
\texttt{--snapshot}& Checkpoint CID             & \texttt{--snapshot cls1-sha256-…} \\
\bottomrule
\end{tabular}
\caption{Glue time flags — enable time-travel queries against the kernel's history}
\end{table}

\subsection{Output Flags}

Output format is controlled by a single flag. Every command produces consistent,
parseable output in all modes:

\begin{table}[H]
\centering
\renewcommand{\arraystretch}{1.35}
\begin{tabular}{@{}lp{10cm}@{}}
\toprule
\textbf{Flag} & \textbf{Output format} \\
\midrule
(default)      & Compact human-readable table \\
\texttt{--summary}   & One-line status line \\
\texttt{--ops}       & Operations-focused: costs, budgets, throughput \\
\texttt{--forensic}  & Full audit trail with CIDs and signatures \\
\texttt{--narrative} & Prose explanation for human readers \\
\texttt{--json}      & Machine-readable JSON for automation \\
\bottomrule
\end{tabular}
\caption{Glue output format flags — same command, different audiences}
\end{table}

\section{connectorctl: The CLI}

\texttt{connectorctl} is the primary operator interface for Connector nodes.
It is a single compiled binary that implements the full Glue grammar.

\begin{lstlisting}[language=bash, caption=connectorctl example session]
# Check node health
connectorctl health

# List all running agents
connectorctl list agents

# Show a specific agent
connectorctl show agent medical-assistant-1

# Trace decisions in the last 15 minutes
connectorctl trace agent medical-assistant-1 --last 15m

# Generate a cryptographic proof
connectorctl prove agent medical-assistant-1 --json

# Explain the last decision made by an agent
connectorctl explain medical-assistant-1

# Deploy an agent from a manifest
connectorctl deploy agent ./manifests/triage-agent.toml

# Inspect with full forensic detail
connectorctl inspect agent medical-assistant-1 --forensic
\end{lstlisting}

\subsection{Standardised Output Contracts}

Every \texttt{connectorctl} command produces output conforming to a stable
contract. Operators can pipe output to \texttt{jq}, monitoring systems, or
alerting pipelines without special-casing each command:

\begin{lstlisting}[language=bash]
# Human-readable table (default)
connectorctl list agents
  ID                      STATUS    UPTIME    COST TODAY
  medical-assistant-1     running   2h 15m    $0.42
  data-pipeline-2         running   5h 42m    $1.18
  audit-agent-3           paused    ---       $0.07

# JSON output for automation
connectorctl list agents --json
{
  "agents": [...],
  "metadata": { "request_id": "req-abc", "duration_ms": 12 }
}
\end{lstlisting}

\section{The REST API Surface}

Every Glue verb is available over HTTP at \texttt{http://localhost:9091/api/v1/}.
The URL structure mirrors the CLI grammar: \texttt{verb/noun} → endpoint path.

\begin{table}[H]
\centering
\renewcommand{\arraystretch}{1.4}
\rowcolors{2}{gray!5}{white}
\begin{tabular}{@{}llp{6cm}@{}}
\toprule
\textbf{Method} & \textbf{Path} & \textbf{Action} \\
\midrule
GET  & \texttt{/agents}                & List all agents \\
POST & \texttt{/agents}                & Register a new agent \\
GET  & \texttt{/agents/:pid}           & Show agent state \\
POST & \texttt{/agents/:pid/start}     & Start agent \\
POST & \texttt{/agents/:pid/chat}      & Invoke agent (LLM gateway) \\
POST & \texttt{/agents/:pid/decisions} & Record a governance decision \\
GET  & \texttt{/agents/:pid/audit/receipts} & Audit receipt chain \\
GET  & \texttt{/agents/:pid/cost}      & Token cost summary \\
GET  & \texttt{/agents/:pid/books}     & Trust score and integrity report \\
GET  & \texttt{/disputes/:id/report}   & Compliance report for a decision \\
GET  & \texttt{/disputes/:id/defense-package} & Full defense package \\
POST & \texttt{/proofs}                & Generate cryptographic proof \\
GET  & \texttt{/actionlog/regulation-report/:reg} & HIPAA/SOC2/GDPR report \\
GET  & \texttt{/health/live}           & Liveness probe \\
GET  & \texttt{/health/ready}          & Readiness probe \\
\bottomrule
\end{tabular}
\caption{Core REST API endpoints — the full Glue surface over HTTP}
\end{table}

All responses follow a consistent envelope:

\begin{lstlisting}[language=bash]
{
  "ok": true,
  "data": { ... },
  "error": null,
  "metadata": { "request_id": "req-xyz", "duration_ms": 8 }
}
\end{lstlisting}

\section{Configuration Hierarchy}

Glue configuration resolves through four layers, with later layers overriding
earlier ones. This means sensible defaults work out of the box, and individual
commands can always be overridden without changing global config:

\begin{figure}[H]
\centering
\begin{tikzpicture}[
  layer/.style={rectangle, rounded corners=3pt, draw=#1!70, fill=#1!8,
                text width=10cm, align=left, minimum height=0.9cm,
                font=\footnotesize, inner sep=7pt},
  node distance=0.4cm
]
\node[layer=CInfra]  (d1) {\textbf{1. Built-in defaults}\enspace — zero-config local development};
\node[layer=CKernel, below=of d1] (d2)
  {\textbf{2. Config file}\enspace \texttt{\textasciitilde/.connector/config.toml} or \texttt{/etc/connector/config.toml}};
\node[layer=CExec, below=of d2] (d3)
  {\textbf{3. Environment variables}\enspace \texttt{CONNECTOR\_API\_URL}, \texttt{CONNECTOR\_API\_KEY}, \texttt{CONNECTOR\_DEV\_MODE}, \ldots};
\node[layer=CDX, below=of d3] (d4)
  {\textbf{4. CLI flags}\enspace \texttt{--api-url}, \texttt{--api-key}, \texttt{--output}, \ldots (highest priority)};

\node[right=0.3cm of d4, font=\tiny\itshape, text=CDX] {highest};
\node[right=0.3cm of d1, font=\tiny\itshape, text=CInfra] {lowest};

\end{tikzpicture}
\caption{Configuration resolution order — later layers override earlier ones}
\label{fig:config-hierarchy}
\end{figure}

\section{Structured Error Model}

Glue errors are always structured and actionable. Every error carries a
machine-readable code, a human-readable message, and an operator hint:

\begin{lstlisting}[language=bash]
{
  "ok": false,
  "error": {
    "code":    "AGENT_NOT_FOUND",
    "message": "Agent 'medical-assistant-1' not found",
    "hint":    "Run 'connectorctl list agents' to see available agents",
    "see_also": ["https://docs.connector.dev/agents/lifecycle"]
  }
}
\end{lstlisting}

Error codes are namespaced (\texttt{AGENT\_*}, \texttt{SESSION\_*},
\texttt{BUDGET\_*}, \texttt{AUTH\_*}) so automated systems can pattern-match
without parsing message text.

